\chapter{Podsumowanie}
\label{chap:podsumowanie}

W ramach projektu zaprojektowano, skonfigurowano oraz zweryfikowano maszynę wirtualną \textbf{VulnBox} opartą na systemie \textbf{Ubuntu Server 26.04.1 LTS}, spełniającą wszystkie wymagania funkcjonalne i bezpieczeństwa.

\section{Macierz zgodności z wymaganiami projektowymi}
\label{sec:macierz-zgodnosci}

Tabela \ref{tab:macierz-zgodnosci} przedstawia zbiorcze podsumowanie weryfikacji poszczególnych wymagań technicznych wraz ze wskazaniem miejsca ich udokumentowania oraz wynikiem testów.

\begin{table}[H]
    \centering
    \caption{Macierz spełnienia wymagań projektowych dla maszyny VulnBox}
    \label{tab:macierz-zgodnosci}
    \tablefont
    \begin{tabularx}{\textwidth}{L{4.2cm} Y C{2.2cm} C{1.6cm}}
        \toprule
        \hcell{Wymaganie projektowe} & \hcell{Sposób weryfikacji technicznej} & \hcell{Rozdział} & \hcell{Status} \\
        \midrule
        Ubuntu Server 26.04.1 LTS bez GUI & Instalacja serwerowa, brak procesów graficznych X11/Wayland & Rozdz. \ref{chap:wstep}, \ref{chap:konfiguracja} & \cmark \\
        Kompatybilność z VirtualBox (\texttt{.ova}) & Eksport OVF 2.0 i testowy re-import pliku \texttt{VulnBox.ova} & Rozdz. \ref{chap:optymalizacja} & \cmark \\
        Automatyczna konfiguracja sieci (DHCP) & Dopasowanie \texttt{en*}/\texttt{eth*} z \texttt{dhcp-identifier: mac} w Netplan & Rozdz. \ref{chap:konfiguracja}, \ref{chap:optymalizacja} & \cmark \\
        Zabezpieczenie programu GRUB2 & Weryfikacja \texttt{password\_pbkdf2} i flagi \texttt{--unrestricted} & Rozdz. \ref{chap:konfiguracja}, \ref{chap:optymalizacja} & \cmark \\
        Aktywna zapora UFW i Fail2Ban & Sprawdzenie reguł TCP 22, 445, 8080 oraz profilu \texttt{sshd} & Rozdz. \ref{chap:konfiguracja}, \ref{chap:optymalizacja} & \cmark \\
        Wieloetapowa ścieżka (\texttt{0-knowledge} $\rightarrow$ \texttt{root}) & Test E2E: Samba $\rightarrow$ Nginx/PHP $\rightarrow$ Flask API $\rightarrow$ \texttt{dev} $\rightarrow$ \texttt{root} & Rozdz. \ref{chap:architektura}, \ref{chap:weryfikacja} & \cmark \\
        Flagi \texttt{user.txt} i \texttt{root.txt} (\texttt{0600}) & Kontrola uprawnień \texttt{stat -c '\%a:\%U'} oraz odczyt w teście E2E & Rozdz. \ref{chap:weryfikacja}, \ref{chap:optymalizacja} & \cmark \\
        Automatyzacja budowy i CI/CD & Skrypty \texttt{provision\_vm.sh}, \texttt{verify\_vm.sh} oraz potoki GitHub Actions & Rozdz. \ref{chap:wstep}, \ref{chap:optymalizacja} & \cmark \\
        \bottomrule
    \end{tabularx}
\end{table}

\section{Wnioski końcowe}
\label{sec:wnioski}

Realizacja projektu potwierdziła, że przygotowanie maszyny wirtualnej wymaga oddzielenia celowo zaprojektowanych słabości aplikacyjnych od zabezpieczeń samego systemu operacyjnego i procesu rozruchowego:
\begin{enumerate}
    \item \textbf{Powtarzalność i automatyzacja wdrożenia:} Oparcie konfiguracji na \textit{Infrastructure as Code} w repozytorium GitHub projektu (\url{https://github.com/alehshehda/APNWA-VulnBox}) \cite{vulnbox_repo}, zawierającym szczegółową instrukcję budowy (\texttt{box/SETUP.md}), przewodnik weryfikacyjny (\texttt{box/WALKTHROUGH.md}) oraz skrypty \texttt{provision\_vm.sh}, \texttt{verify\_vm.sh} i \texttt{cleanup\_vm.sh}, pozwala na odtworzenie identycznego obrazu maszyny w dowolnym momencie oraz automatyczną walidację usług w potoku CI/CD (\texttt{box-ci.yml}).
    \item \textbf{Realizm ścieżki weryfikacji:} Połączenie czterech odmiennych technologii (udział plikowy SMB, aplikacja PHP z błędem normalizacji ścieżek CWE-22, odizolowany na interfejsie \texttt{127.0.0.1} serwis Flask oraz błędy uprawnień w zadaniach \texttt{cron} i regułach \texttt{sudoers.d}) wymusza przeprowadzenie pełnego procesu analitycznego od rekonesansu po weryfikację uprawnień lokalnych.
    \item \textbf{Odporność obrazu dystrybucyjnego (\texttt{VulnBox.ova}):} Zastosowanie uniwersalnych reguł Netplan niewrażliwych na zmianę adresu MAC oraz zabezpieczenie GRUB2 skrótem PBKDF2 gwarantują, że maszyna uruchamia się bezobsługowo po imporcie w programie VirtualBox, nie pozwalając przy tym na obejście ścieżki przez edycję parametrów startowych jądra.
\end{enumerate}
